%\documentclass[12pt,oneside]{amsart}
\documentclass{scrartcl}

%\documentclass{article}

\usepackage{amsthm,amsmath,amssymb}
\usepackage[numbers]{natbib}

%\usepackage[colorlinks]{hyperref}
\usepackage{hyperref}
\hypersetup{
    colorlinks,%
    citecolor=black,%
    filecolor=black,%
    linkcolor=black,%
    urlcolor=black
}
\usepackage{hypernat}
\usepackage[top=2.5cm, bottom=2.5cm, left=2.8cm,
right=2.8cm]{geometry}
\usepackage{graphicx}


%\setlength{\parskip}{0pt}
%\setlength{\parsep}{0pt}
%\setlength{\headsep}{0pt}
%\setlength{\topskip}{0pt}
%\setlength{\topmargin}{0pt}
%\setlength{\topsep}{0pt}
%\setlength{\partopsep}{0pt}

\linespread{1}

%\usepackage{marvosym}

%\textwidth = 6.5 in \textheight = 9.2 in \oddsidemargin = 0.0 in
%\evensidemargin = 0.0 in \topmargin = -0.0 in \headheight = 0.0 in
%\headsep = 0.0in \parskip = 0.0in \parindent = 0.0in
%\def\baselinestretch{0.871}


\newtheorem{theorem}{Theorem}[section]
\newtheorem{proposition}[theorem]{Proposition}
\newtheorem{problem}[theorem]{Problem}
\newtheorem{fact}[theorem]{Fact}
\newtheorem{lemma}[theorem]{Lemma}
\newtheorem{defn}[theorem]{Definition}
\newtheorem{corollary}[theorem]{Corollary}
\newtheorem{remark}{Remark}[section]
\newtheorem{example}[theorem]{Example}

\newcommand{\E}{{\mathbb E}}
\newcommand{\se}{{\mathcal{E}}}
\newcommand{\R}{{\mathbb R}}
\renewcommand{\P}{{\mathbb P}}
\newcommand{\ac}{{\mathcal{A}}}
\newcommand{\A}{{\mathcal{A}}}
\newcommand{\B}{{\mathcal{B}}}
\newcommand{\C}{{\mathcal{C}}}
\newcommand{\M}{{\mathcal{M}}}
\newcommand{\Mf}{{\mathfrak{M}}}
\newcommand{\T}{{\mathcal{T}}}
\newcommand{\Z}{{\mathcal{Z}}}
\newcommand{\K}{{\mathcal{K}}}
\newcommand{\lc}{{\mathcal{L}}}
\newcommand{\Ps}{{\mathcal{P}}}
\newcommand{\G}{{\mathcal{G}}}
\newcommand{\pa}{{\mathring{p}}}
\newcommand{\F}{{\cal F}}
\newcommand{\samp}{{\mathcal{X}}}
\newcommand{\X}{{\mathcal{X}}}
\newcommand{\hi}{{\hat{\Phi}}}
\newcommand{\data}{{\text{data}}}

\newcounter{rcnt}[section]
\renewcommand{\thercnt}{(\roman{rcnt})}

\def\qt#1{\qquad\text{#1}}

\def\argmin{\mathop{\rm argmin}}
\def\argmax{\mathop{\rm argmax}}


\setlength{\parskip}{1.4 \medskipamount}
%\sloppy
%\linespread{1.3}

\begin{document}

\title{STAT 238 - Bayesian Statistics  \\
  Lecture Twelve} 
\subtitle{Spring 2026, UC Berkeley}

\author{Aditya Guntuboyina}


\date{18 Feb 2026}

\maketitle

Today we shall generalize the Beta-Binomial analysis from last week to
Dirichlet-Multinomial inference. First we shall look at the
Multinomial distribution (which is a generalization of the Binomial
distribution) and the Dirichlet distribution (which is a
generalization of the Beta distribution).

\section{Multinomial Distribution}

The multinomial distribution is a generalization of the Binomial
distribution. We first recall the binomial distribution. We say that
$X \sim \text{Bin}(n, p)$ if
\begin{align*}
   \P\{X = x\} = \frac{n!}{x! (n-x)!} p^x (1 - p)^{n-x} ~~ \text{ for
  } x = 0, 1, \dots, n. 
\end{align*}
$X$ represents the count of successess in $n$ repetitions of a simple
binary experiment with two outcomes success (probability $p$) and
failure (probability $1 - p$).

The multinomial distribution is obtained by considering $n$
repetitions of an experiment with $k$ outcomes (for some $k \geq 1$)
with probabilities $p_1, \dots, p_k$ (we need $p_i \geq 0$ and
$\sum_{i=1}^kp_i = 1$). Let $X_i$ denote the count of the $i$-th
outcome among the $n$ repetitions of the experiment (in other words,
$X_i$ is the number of times the $i$-th outcome happens in the $n$
trials). The joint distribution of $(X_1, \dots, X_k)$ is written as
$\text{Multinomial}(n; p_1, \dots, p_k)$. It corresponds to the
probabilities:
\begin{align*}
\P\!\left\{ X_1 = x_1, \ldots, X_k = x_k \right\}
&=
\begin{cases}
\displaystyle
\frac{n!}{x_1! \cdots x_k!}
\prod_{i=1}^k p_i^{x_i},
& \text{if } x_i \ge 0 \text{ for all } i
\text{ and } \sum_{i=1}^k x_i = n, \\
0, & \text{otherwise}.
\end{cases}
\end{align*}
The following consequences of $(X_1, \dots, X_k) \sim
\text{Multinomial}(n; p_1, \dots, p_k)$ are straightforward to see:
\begin{enumerate}
\item To derive the distribution of $X_i$ i.e., $\P\{X_i = x\}$, we
  can write
  \begin{align*}
    \P\{X_i = x\} &= \P\{x \text{ trials had outcome } i, (n-x)
                    \text{ trials had outcome not } i \} \\
    &= \frac{n}{x! (n-x)!} p_i^{x} (1 - p_i)^{n-x}. 
  \end{align*}
  In other words $X_i \sim \text{Bin}(n, p_i)$. In essence, here we
  consider the variant of the original experiment where, for each 
  trial, we just record if the original outcome was 'not
  $i$' or $i$. This converts the setup to binary trials and $X_i$
  represents the number of successes (success = outcome is $i$) so
  $X_i \sim \text{Bin}(n, p_i)$. 
\item The joint distribution of $(X_i, X_j)$ (for fixed $i \neq j$)
  can be derived in the following way:
  \begin{align*}
   & \P\{X_i = x_1, X_j = x_2\} \\ &= \P\{x_1 \text{ trials were } i, x_2
                    \text{ trials were } j, n-x_1-x_2 \text{
                                 trials were neither } i \text{
                                     nor } j \}\\
    &= \frac{n!}{x_1! x_2!(n - x_1 - x_2)!} p_i^{x_1} p_j^{x_2} (1 -
      p_i - p_j)^{n-x_1-x_2}
  \end{align*}
  provided $x_1, x_2 \in \{0, 1, \dots, n\}$ with $x_1 + x_2 \leq n$
  (otherwise, the probability equals 0).
\item It can be checked (see e.g.,
  \url{https://en.wikipedia.org/wiki/Multinomial_distribution}) that 
\begin{align*}
  \E X_i = n p_i, ~~ \text{var}(X_i) = np_i(1 - p_i), \text{ and } ~~ \text{cov}(X_i,
  X_j) = -np_i p_j \text{ if } i \neq j. 
\end{align*}
\end{enumerate}

\section{Dirichlet Distribution}
The Dirichlet distribution is a generalization of the Beta
distribution. Recall the Beta$(a, b)$ distribution is a distribution
over $p \in [0, 1]$ corresponding to the density:
\begin{align*}
  f(p) = \frac{\Gamma(a + b)}{\Gamma(a) \Gamma(b)} p^{a-1} (1 -
  p)^{b-1} I\{0 \leq p \leq 1\}. 
\end{align*}
This density function is well-defined only when both $a$ and $b$ are
strictly positive. However one can still define the Beta$(a, b)$
distribution even when one or both of $a$ and $b$ are zero. This can
be done by a limiting argument: 
\begin{enumerate}
\item $\text{Beta}(0, 0) = \lim_{\epsilon \downarrow 0}
  \text{Beta}(\epsilon, \epsilon) = \text{Bernoulli}(0.5)$. In other
  words, $\text{Beta}(0, 0)$ becomes a discrete two-point distribution assigning
  equal probability to 0 and 1.
\item For fixed $b > 0$, we have $\text{Beta}(0, b) = \lim_{\epsilon
    \downarrow 0} \text{Beta}(\epsilon, b) = \text{Bernoulli}(0) = 
  \delta_{\{0\}}$. In other words, $\text{Beta}(0, b)$ is the point
  mass as 0.
\item For fixed $a > 0$, we have $\text{Beta}(a, 0) = \lim_{\epsilon
    \downarrow 0} \text{Beta}(a, \epsilon) = \text{Bernoulli}(1) = 
  \delta_{\{1\}}$. In other words, $\text{Beta}(a, 0)$ is the point
  mass as 1.  
\end{enumerate}
The limiting statements above can be made rigorous by
moment calculations (e.g., by showing that the moments of
$\text{Beta}(\epsilon, \epsilon)$ converge to the moments of
$\text{Bernoulli}(0.5)$ as $\epsilon \downarrow 0$).

The Beta distribution can be viewed as the distribution over the
probabilities ($p$ and $1-p$) of an experiment with two outcomes. The Dirichlet
distribution is the distribution over the probabilities $p_1, \dots,
p_k$ of an experiment with $k$ outcomes. Specifically given $a_1,
\dots, a_k$, the Dirichlet$(a_1, \dots, a_k)$ distribution corresponds
to the density:
\begin{align*}
  \frac{\Gamma(a_1 + \dots + a_k)}{\Gamma(a_1) \dots \Gamma(a_k)}
  p_1^{a_1 - 1} \dots p_k^{a_k - 1} I\{p_1, \dots, p_k \geq 0,
  \sum_{i} p_i = 1\}.  
\end{align*}
Note that this should be viewed as density on $(k-1)$-dimensional
space (as opposed to $k$-dimensional space) because of the restriction
$p_1 + \dots p_k = 1$. In other words, for every $A \subseteq \R^k$,
we have
\begin{align*}
  &  \P\{(p_1, \dots, p_k) \in A\} \\
  &=  \frac{\Gamma(a_1 + \dots + a_k)}{\Gamma(a_1) \dots \Gamma(a_k)} \int_S  
  p_1^{a_1 - 1} \dots p_{k-1}^{a_{k-1} - 1} \left(1 - p_1 - \dots -
    p_{k-1} \right)^{a_k - 1} \\ &I\{(p_1, \dots, p_{k-1}, 1 - p_1 - \dots -
                                   p_{k-1}) \in A\} dp_1 \dots 
    dp_{k-1}
\end{align*}
where
\begin{align*}
  S := \left\{(p_1, \dots, p_{k-1}) : p_{i} \geq 0, p_1 + \dots +
  p_{k-1} \leq 1 \right\}. 
\end{align*}
It can be checked that if $(p_1, \dots, p_k) \sim
\text{Dirichlet}(a_1, \dots, a_k)$, then (see e.g.,
\url{https://en.wikipedia.org/wiki/Dirichlet_distribution}) 
\begin{align*}
  &\E p_i = \frac{a_i}{a_1 + \dots + a_k} \\ & \text{var}(p_i) =
                                               \left(\frac{a_i}{a_1 + \dots
                                               + a_k} \right) \left( \frac{(a_1 +
                                               \dots + a_k) - a_i}{a_1
                                               + \dots + a_k} \right) \left(
                                               \frac{1}{a_1 + \dots +
                                               a_k + 1} \right) \\
  &\text{cov}(p_i, p_j) = -\left(\frac{a_i}{a_1 + \dots + a_k} \right) \left(
    \frac{a_j}{a_1 + \dots + a_k} \right) \left( \frac{1}{a_1 + \dots
    + a_k + 1} \right). 
\end{align*}


As for the case of the Beta, the following facts about the Dirichlet
can also be proved using moment calculations:
\begin{enumerate}
\item $\text{Dirichlet}(0, \dots, 0) = \lim_{\epsilon \downarrow 0}
  \text{Dirichlet}(\epsilon, \dots, \epsilon)$ is the discrete uniform
  distribution on $e_1 \dots, e_k$ with $e_i$ is the vector which has
  $1$ in the $i$-th location and 0 everywhere else i.e.,
  \begin{align*}
    \text{Dirichlet}(0, \dots, 0) = \frac{1}{k} \sum_{i=1}^k
    \delta_{\{e_i\}}. 
  \end{align*}
\item If $k = 4$ and $a_2, a_4 > 0$, then $\text{Dirichlet}(0, a_2, 0,
  a_4) = \lim_{\epsilon \downarrow 0} \text{Dirichlet}(\epsilon, a_2,
  \epsilon, a_4)$ is supported on $\{(p_1, p_2, p_3, p_4) : p_1 = 0, p_2 \geq 0,
  p_3 = 0, p_4 \geq 0, p_1 + p_2 + p_3 + p_4 = 1 \}$, and the
  distribution of $(p_2, p_4)$ is $\text{Dirichlet}(a_2,
  a_4)$. Informally, we write
  \begin{align*}
    \text{Dirichlet}(0, a_2, 0, a_4) = \text{Dirichlet}(a_2, a_4). 
  \end{align*}
\item More generally, let $a = (a_1, \dots, a_k)$  with each $a_i \geq
  0$. Let $I = \{i: a_i > 0\}$ and let $m$ be the cardinality of
  $I$. Then $\text{Dirichlet}(a) := \lim_{\epsilon \downarrow 0}
  \text{Dirichlet}(a_1 + \epsilon, \dots, a_k + \epsilon)$ satisfies
  the following. 
  \begin{enumerate}
  \item The support of $\text{Dirichlet}(a)$ equals $\{(p_1, \dots,
    p_k): p_i \geq 0, \sum_{i=1}^k p_i = 1, p_i = 0 \text{ for all } i
    \notin I\}$.
  \item If $p \sim \text{Dirichlet}(a)$, then the subvector $p_i, i
    \in I$ satisfies:
    \begin{align*}
      (p_i, i \in I) \sim \text{Dirichlet}(a_i, i \in S). 
    \end{align*}
    Note that the above is a Dirichlet distribution for an
    $m$-dimensional probability vector. 
  \end{enumerate}
\end{enumerate}

\section{Dirichlet Prior and Multinomial Likelihood}

We saw the following basic fact in Lecture 8:
\begin{align*}
  p \sim \text{Beta}(a, b) ~~ \text{ and } ~~ X \mid p \sim
  \text{Bin}(n, p) \implies p \mid X = x \sim \text{Beta}(a + x, b + n
  - x). 
\end{align*}
The generalization of this is:
\begin{align*}
&  (p_1, \dots, p_k) \sim \text{Dirichlet}(a_1, \dots, a_k)  \text{
  and }  (X_1, \dots, X_k) \mid (p_1, \dots, p_k) \sim 
  \text{Multinomial}(n; p_1, \dots, p_k) \\
& \implies (p_1, \dots, p_k) \mid X_1 = x_1, \dots, X_k = x_k \sim
  \text{Dirichlet}(a_1 + x_1, a_2 + x_2, \dots, a_k + x_k). 
\end{align*}
So the posterior mean estimate of $p_i$ is given by:
\begin{align*}
  \E (p_i \mid X_1 = x_1, \dots, X_k = x_k) = \frac{a_i + x_i}{(a_1 +
  x_1) + \dots + (a_k + x_k)} = \frac{a_i + x_i}{n + a_1 + \dots +
  a_k}
\end{align*}
where we used $x_1 + \dots + x_k = n$.

In the special case where all $a_i = 0$, then the posterior mean of
$p_i$ is simply $x_i/n$; so we are estimating the probability of the
$i$-th outcome by simply the proportion of times the $i$-th outcome
appeared in the $n$ trials. This is just the frequentist MLE. So the
posterior mean estimate coincides with the frequentist MLE when the
prior is $\text{Dirichlet}(0, \dots, 0)$.

More precisely, the posterior corresponding to the
$\text{Dirichlet}(0, \dots, 0)$ prior equals  $\text{Dirichlet}(x_1,
\dots, x_k)$. A key property of this posterior is that if $x_i = 0$ for
some $i$, then, by the properties of Dirichlet distributions with some
zero hyperparameters discussed in the previous section, the posterior
places all its mass on the set $\{p_i = 0\}$. In other words, any
outcome that is not observed in the sample is assigned zero
probability under the posterior and is therefore completely excluded
from future consideration.





%\begin{thebibliography}{99}

  
%\bibitem[Fisher(1934)]{Fisher1934}
%R.~A. Fisher,
%\newblock ``Probability, likelihood and quantity of information in the logic of
%uncertain inference,''
%\newblock \emph{Proceedings of the Royal Society of London. Series~A,
%Containing Papers of a Mathematical and Physical Character}, vol.~146,
%no.~856, pp.~1--8, 1934.

%\bibitem[Efron(2010)]{Efron}
%Efron, B. (2010)
%\textit{Large-scale inference: empirical Bayes methods for estimation,
%testing, and prediction}. Cambridge University Press, 2010.
  
%     \bibitem[Jaynes(2003)]{Jaynes} Jaynes, E. T, (2003)
%  \textit{Probability theory: the logic of science}. Cambridge
%  University Press, 2003.

%  \bibitem[Sinai(1992)]{Sinai} Sinai, Y. G, (1992)
%  \textit{Probability theory: an introductory course}. Springer
%  Verlag, 1992.  

%\bibitem[{deGroot(1986)}]{DeGrootBlackwell}DeGroot, M. H. (1986)
%       A conversation with David Blackwell.
%\newblock \emph{Statistical Science}, \textbf{1}, 40--53.

%         \bibitem[Mosteller(1987)]{Mosteller} Mosteller, F, (1987)
%  \textit{Fifty challenging problems in probability with
%    solutions}. Courier Corporation, 1987.

%    \bibitem[MacKay(2003)]{MacKay} MacKay, D. J. C., (2003)
%  \textit{Information theory, inference, and learning
%  algorithms}. Cambridge University Press, 2003.

%\bibitem[{Lindley(2013)}]{Lindley}Lindley, D. V. (2013)
%   \textit{Understanding Uncertainty}. John Wiley and sons, 2013.


%\end{thebibliography}







\end{document}

%%% Local Variables:
%%% mode: latex
%%% TeX-master: t
%%% End:
